% Generated draft: requires mathematical and source review.
\documentclass[11pt]{article}
\usepackage[margin=25mm]{geometry}
\usepackage{fontspec}
\setmainfont{Latin Modern Roman}
\usepackage{amsmath,amssymb,mathtools,unicode-math}
\setmathfont{Latin Modern Math}
\usepackage{enumitem,needspace,xurl,hyperref}
\setlength{\parindent}{0pt}
\setlength{\parskip}{5pt}
\setlength{\emergencystretch}{3em}
\newcommand{\N}{\mathbb{N}}
\newcommand{\Z}{\mathbb{Z}}
\newcommand{\Q}{\mathbb{Q}}
\newcommand{\R}{\mathbb{R}}
\newcommand{\C}{\mathbb{C}}
\newcommand{\F}{\mathbb{F}}
\newcommand{\A}{\mathbb{A}}
\newcommand{\PP}{\mathbb{P}}
\newcommand{\E}{\mathbb{E}}
\newcommand{\eps}{\varepsilon}
\newcommand{\dd}{\,\mathrm{d}}
\begin{document}
\section*{UnsolvedMath ID: 30004271}
\textbf{Source problem code:} catalogue entry
\subsection{Definitions and mathematical statement}
Let \(M\) carry an almost-Riemannian structure. In the usual local description, this structure is represented by \(n=\dim M\) smooth vector fields \(X_1,\ldots,X_n\), declared orthonormal wherever they are linearly independent. Its regular region is
\[
R=\{p\in M:\dim\operatorname{span}\{X_1(p),\ldots,X_n(p)\}=n\}.
\]
On \(R\), the structure induces a Riemannian metric \(g\). Write \(\mu_g\) for its Riemannian volume measure and \(L^2(R,\mu_g)\) for the Hilbert space of square-integrable complex-valued functions. Let \(C_c^\infty(R)\) denote the smooth functions with compact support in \(R\). With the sign convention \(\Delta_g=\operatorname{div}_g\operatorname{grad}_g\), consider the densely defined symmetric operator
\[
A:C_c^\infty(R)\subset L^2(R,\mu_g)\longrightarrow L^2(R,\mu_g),
\qquad Au=\Delta_g u.
\]
Essential self-adjointness means that the operator closure \(\overline A\) is self-adjoint, equivalently \(\overline A=A^*\). Reversing the sign of the Laplacian does not change this property.

\textbf{Question as supplied.} For every almost-Riemannian structure in the intended class, is \(A\) essentially self-adjoint? In particular, does this hold for every real-analytic almost-Riemannian structure on a two-dimensional manifold having no tangency points?

For the usual two-dimensional situation in which the singular set \(\Sigma=M\setminus R\) is locally a smooth curve and the distribution
\[
\mathcal D_p=\operatorname{span}\{X_1(p),X_2(p)\}
\]
has dimension one at \(p\in\Sigma\), a tangency point is a point satisfying \(\mathcal D_p=T_p\Sigma\). Thus absence of tangency points means \(\mathcal D_p\ne T_p\Sigma\) at every singular point to which this description applies.

\textbf{Unresolved scope qualifications.} The supplied record does not specify the complete definition of the admissible almost-Riemannian structures, including bracket-generation or singular-set regularity hypotheses. The description of tangency above therefore explains the usual smooth-curve case without resolving the treatment of more degenerate singularities. No compactness, completeness, boundary, or control-of-ends assumptions are supplied. These omissions prevent identification of a fully specified universal class and must be checked against the source before treating the question as a precise conjecture. No such hypotheses are added here, and no boundary conditions or self-adjoint extension are selected.
\subsection{Short English statement}
The question asks whether the Laplace--Beltrami operator on the regular region of an almost-Riemannian structure has a unique self-adjoint extension from compactly supported smooth functions. It singles out real-analytic surfaces without tangency points. The precise admissible structures and global hypotheses remain unspecified. The supplied \emph{partially solved} label and literature assessment are dated curation, not independently verified conclusions.
\subsection{Sources}
\begin{itemize}
\item Local UnsolvedMath registry: supplied complete record with permanent ID \(30004271\), source problem code \texttt{catalogue entry}, title \emph{Essential Self-Adjointness on Almost-Riemannian Regular Regions}, in the Oberwolfach Reports Open Problems collection. The record gives proposed year 2019 and names no individual proposer.
\item Supplied report citation: \emph{Mini-Workshop: Self-adjoint Extensions in New Settings} (2020), Oberwolfach Reports, pp.~2911--2950, DOI \texttt{10.4171/owr/2019/47}. \url{https://doi.org/10.4171/owr/2019/47}. The supplied background describes Oberwolfach Reports as documenting mathematical workshops and their open-problem discussions.
\item Supplied literature source: Franceschi--Prandi--Rizzi. \url{https://doi.org/10.1007/s11118-018-09760-w}.
\item Supplied literature source: OWR 2019 47. \url{https://ems.press/content/serial-article-files/46828}.
\item Supplied literature source: \emph{Quantum confinement on non-complete Riemannian manifolds}. \url{https://ems.press/content/serial-article-files/33671?nt=1}.
\item Supplied literature assessment, checked \textbf{2026-08-22}: the record labels the problem \emph{partially solved} and reports essential self-adjointness under important regularity and non-characteristic hypotheses, while describing the broad real-analytic no-tangency conjecture in the extracted record as unresolved for general almost-Riemannian structures. This is supplied dated curation derived from the OWR material, not an independently verified literature assessment. The supplied background expressly requests independent verification of time-sensitive status and cited claims.
\item The registry labels its statement \texttt{corrected\_verified}, with an explanation that mathematical content was preserved while notation, typography, duplicate clauses, or omitted definitions were repaired. This is a supplied editorial qualification, not mathematical verification of this draft or of the conjecture.
\end{itemize}
\end{document}
